\uuid{vQKw}
\titre{Propriétés de la fonction de répartition de la loi normale centrée réduite}
\auteur{Maxime Nguyen}
\organisation{AMSCC}
\datecreate{2026-11-27}
\niveau{L2}
\module{Probabilité}
\chapitre{Probabilité continue}
\sousChapitre{Loi normale}
\difficulte{2}
\theme{loi normale centrée réduite, fonction de répartition, monotonie}

\contenu{
\begin{question}
Soit $\varphi(x) = \frac{1}{\sqrt{2\pi}}\,e^{-\frac{x^2}{2}}$ la densité de la loi normale centrée réduite $\mathcal{N}(0,1)$ et $\Phi$ sa fonction de répartition. Quelles affirmations sont vraies ?
\begin{answers}
\good{$\Phi$ est strictement croissante sur $\mathbb{R}$.}
\bad{Il existe un réel $t$ tel que $\Phi(t) = 1$.}
\good{$\Phi$ est dérivable sur $\mathbb{R}$ et $\Phi' = \varphi$.}
\bad{$\Phi$ est concave sur $\mathbb{R}$.}
\bad{Pour tout réel $t$, $\Phi(t) = \displaystyle\int_0^t \varphi(x)\,\mathrm{d}x$.}
\end{answers}
\begin{explanations}
A, C : $\varphi$ est continue et strictement positive sur $\mathbb{R}$, donc $\Phi$ est dérivable de dérivée $\varphi > 0$.

B est faux : $1 - \Phi(t) = \int_t^{+\infty}\varphi > 0$ pour tout $t$. $\Phi(t)$ tend vers $1$ sans l'atteindre ; une table qui affiche $1{,}0000$ arrondit.

D est faux : $\Phi'' = \varphi'$, et $\varphi'(x) = -x\,\varphi(x)$ est positive pour $x < 0$. $\Phi$ est convexe sur $]-\infty;0]$ et concave sur $[0;+\infty[$.

E est faux : $\Phi(t) = \int_{-\infty}^t \varphi$ ; l'intégrale de $0$ à $t$ vaut $\Phi(t) - \frac12$.

Contre-sens visés : croire qu'une probabilité de queue devient nulle « loin »
(B) ; oublier la borne $-\infty$ dans la définition de la fonction de
répartition (E).
\end{explanations}
\end{question}
}
